\renewcommand{\labelenumii}{\arabic{enumi}.\arabic{enumii}}
\renewcommand{\labelenumiii}{\arabic{enumi}.\arabic{enumii}.\arabic{enumiii}}
\renewcommand{\labelenumiv}{\arabic{enumi}.\arabic{enumii}.\arabic{enumiii}.\arabic{enumiv}}

\mychapter{Anticipated Outputs}{Anticipated Outputs}{}
\label{chap:Anticipated Outputs}

\mysection{Anticipated Outputs}{Anticipated Outputs}
\label{sec:background}

This research project is structured to deliver contributions to academic dissemination, software infrastructure and practical implementation guidelines for the biomechanics community
\subsection{Academic Publications}

\subsubsection{Conference Proceedings}

\textbf{South African Society of Biomechanics Conference (SASB 2025)}

\textbf{REAL-TIME 3D BIOMECHANICAL ANALYSIS AND AI PERSONAL TRAINER FOR STRENGTH AND
CONDITIONING EXERCISES USING MARKERLESS MOTION CAPTURE.
}
We contributed to the aforementioned conference paper that presented a novel, real-time markerless
3D motion capture system that provides biomechanical
feedback using low-cost hardware and open-source
software. Whilst initial results are promising, limitations such
as sensitivity to video quality and the need for
broader validation remain. 
We are looking to address these limitations by the continuation of this research with hopes of improving the quality of the video input by identifying optimal camera configurations.


\subsubsection{Journal Articles}

This work presents the novel combination of B2 and Simulated Annealing to solve the camera placement problem. A potential title is \textit{Synthetic annealing: A simulation-based framework for optimizing multi-view camera configurations in markerless motion capture.}

A key contribution is proving that there are prefered camera setups for markerless systems. This would allow other markerless systems to improve by optimally configuring their camera network.

\subsubsection{Potential Avenues for future publications}

Beyond the primary methodological paper, the modular nature of the proposed simulation pipeline enables several derivative studies. These include the following contributions

\begin{enumerate}
    \item A biomechanical analysis of exercise-specific occlusion patterns to derive task-specific camera setups.
    \item A comparative study evaluating how different state-of-the-art pose estimation algorithms degrade under varying camera angles.
\end{enumerate}

Collectively, these outputs represent a significant transition from heuristic-based sensor deployment to mathematically optimized configuration strategies. By quantifying the relationship between camera placement and reconstruction accuracy, this research aims to democratize access to high-fidelity motion capture in resource-constrained environments. The synthesis of these academic, technical, and practical contributions will culminate in the final MEng thesis submission, titled "Improving markerless pose estimation systems by identifying optimal camera configurations," scheduled to be done before November 2026.